\scanpage{129}
% line head | text | 78
\mspage{78}
% line 1–2 | text | Proposition 6
\uline{Proposition 6} L'axiome $W(7)$ (lemme des cinq)
implique $W(7\mathrm{bis})$ (axiome du carré cocartésien).

% line 3 | text | On regarde le diagramme
On regarde le diagramme cubique suivant
% line 4–9 | diagram | cube from intersection, opens and X to repeated U prime and X
\diagram{\begin{tikzcd}[ampersand replacement=\&,row sep=large,column sep=small]
U\cap U'\arrow[rr,"i_0^{\deleted{\unsure{\prime}}}"]\arrow[dr,"i'_0","\deleted{\unsure{i_0}}"']\arrow[dd,"f_{U\cap U'}=i_0"'] \& \& U'\arrow[dr,"i'"]\arrow[dd,"f_{U'}=\mathrm{id}" description] \&\\
\& U\arrow[rr,"i"]\arrow[dd,"f_U=i"'] \& \& X\arrow[dd,"f=\mathrm{id}_X"]\\
U'\arrow[rr,"\mathrm{id}"]\arrow[dr,"i'"'] \& \& U'\arrow[dr,"i'"] \&\\
\& X\arrow[rr,"\mathrm{id}"] \& \& X
\end{tikzcd}}
% line 10–16 | text | du carré cocartésien donné
, du carré cocartésien donné dans le carré cocartésien
évident de $X$ et de ses deux ouverts $X,U'$,
d'intersection $U'$. Ici les deux flèches verticales
$f$ et $f_{U'}$ sont des iso donc \deleted{sont} $\in W$. Avec les
notations de $W(7)$, on a donc c) ($f\in W$ \unsure{!}),
donc a) et b) sont équivalentes, i.e. $i_0\in W$
et $\unsure{\{|(}i\in W$ sont équivalentes \hfill \unsure{qed.}

% line 17–18 | text | Théorème 2
\uline{Théorème 2} Le localiseur fond. $W_{\infty}$ satisfait à
l'axiome $W(7)$.

% line 19–24 | text | C'est sûrement une conséquence
C'est \inserted{sûrement} une conséquence \inserted{du critère \unsure{cohomologique|homotopique} d'Artin-Mazur, compte tenu} \unsure{des|de} la suite exacte
de \unsure{Mayer|Meyer}-Vietoris pour l'un \deleted{\unsure{deux|des}} et l'autre
des deux carrés cocartésiens, et les variantes
non commutatives de \unsure{Meyer|Mayer}-Vietoris. Je n'ai
pas fait la vérification en forme -- et devrais
sans doute réparer cette omission.
% archive sheet number: 129
